% Options for packages loaded elsewhere
% Options for packages loaded elsewhere
\PassOptionsToPackage{unicode}{hyperref}
\PassOptionsToPackage{hyphens}{url}
\PassOptionsToPackage{dvipsnames,svgnames,x11names}{xcolor}
%
\documentclass[
  letterpaper,
  DIV=11,
  numbers=noendperiod]{scrartcl}
\usepackage{xcolor}
\usepackage{amsmath,amssymb}
\setcounter{secnumdepth}{5}
\usepackage{iftex}
\ifPDFTeX
  \usepackage[T1]{fontenc}
  \usepackage[utf8]{inputenc}
  \usepackage{textcomp} % provide euro and other symbols
\else % if luatex or xetex
  \usepackage{unicode-math} % this also loads fontspec
  \defaultfontfeatures{Scale=MatchLowercase}
  \defaultfontfeatures[\rmfamily]{Ligatures=TeX,Scale=1}
\fi
\usepackage{lmodern}
\ifPDFTeX\else
  % xetex/luatex font selection
\fi
% Use upquote if available, for straight quotes in verbatim environments
\IfFileExists{upquote.sty}{\usepackage{upquote}}{}
\IfFileExists{microtype.sty}{% use microtype if available
  \usepackage[]{microtype}
  \UseMicrotypeSet[protrusion]{basicmath} % disable protrusion for tt fonts
}{}
\makeatletter
\@ifundefined{KOMAClassName}{% if non-KOMA class
  \IfFileExists{parskip.sty}{%
    \usepackage{parskip}
  }{% else
    \setlength{\parindent}{0pt}
    \setlength{\parskip}{6pt plus 2pt minus 1pt}}
}{% if KOMA class
  \KOMAoptions{parskip=half}}
\makeatother
% Make \paragraph and \subparagraph free-standing
\makeatletter
\ifx\paragraph\undefined\else
  \let\oldparagraph\paragraph
  \renewcommand{\paragraph}{
    \@ifstar
      \xxxParagraphStar
      \xxxParagraphNoStar
  }
  \newcommand{\xxxParagraphStar}[1]{\oldparagraph*{#1}\mbox{}}
  \newcommand{\xxxParagraphNoStar}[1]{\oldparagraph{#1}\mbox{}}
\fi
\ifx\subparagraph\undefined\else
  \let\oldsubparagraph\subparagraph
  \renewcommand{\subparagraph}{
    \@ifstar
      \xxxSubParagraphStar
      \xxxSubParagraphNoStar
  }
  \newcommand{\xxxSubParagraphStar}[1]{\oldsubparagraph*{#1}\mbox{}}
  \newcommand{\xxxSubParagraphNoStar}[1]{\oldsubparagraph{#1}\mbox{}}
\fi
\makeatother


\usepackage{longtable,booktabs,array}
\usepackage{calc} % for calculating minipage widths
% Correct order of tables after \paragraph or \subparagraph
\usepackage{etoolbox}
\makeatletter
\patchcmd\longtable{\par}{\if@noskipsec\mbox{}\fi\par}{}{}
\makeatother
% Allow footnotes in longtable head/foot
\IfFileExists{footnotehyper.sty}{\usepackage{footnotehyper}}{\usepackage{footnote}}
\makesavenoteenv{longtable}
\usepackage{graphicx}
\makeatletter
\newsavebox\pandoc@box
\newcommand*\pandocbounded[1]{% scales image to fit in text height/width
  \sbox\pandoc@box{#1}%
  \Gscale@div\@tempa{\textheight}{\dimexpr\ht\pandoc@box+\dp\pandoc@box\relax}%
  \Gscale@div\@tempb{\linewidth}{\wd\pandoc@box}%
  \ifdim\@tempb\p@<\@tempa\p@\let\@tempa\@tempb\fi% select the smaller of both
  \ifdim\@tempa\p@<\p@\scalebox{\@tempa}{\usebox\pandoc@box}%
  \else\usebox{\pandoc@box}%
  \fi%
}
% Set default figure placement to htbp
\def\fps@figure{htbp}
\makeatother





\setlength{\emergencystretch}{3em} % prevent overfull lines

\providecommand{\tightlist}{%
  \setlength{\itemsep}{0pt}\setlength{\parskip}{0pt}}



 


\KOMAoption{captions}{tableheading}
\makeatletter
\@ifpackageloaded{caption}{}{\usepackage{caption}}
\AtBeginDocument{%
\ifdefined\contentsname
  \renewcommand*\contentsname{Table of contents}
\else
  \newcommand\contentsname{Table of contents}
\fi
\ifdefined\listfigurename
  \renewcommand*\listfigurename{List of Figures}
\else
  \newcommand\listfigurename{List of Figures}
\fi
\ifdefined\listtablename
  \renewcommand*\listtablename{List of Tables}
\else
  \newcommand\listtablename{List of Tables}
\fi
\ifdefined\figurename
  \renewcommand*\figurename{Figure}
\else
  \newcommand\figurename{Figure}
\fi
\ifdefined\tablename
  \renewcommand*\tablename{Table}
\else
  \newcommand\tablename{Table}
\fi
}
\@ifpackageloaded{float}{}{\usepackage{float}}
\floatstyle{ruled}
\@ifundefined{c@chapter}{\newfloat{codelisting}{h}{lop}}{\newfloat{codelisting}{h}{lop}[chapter]}
\floatname{codelisting}{Listing}
\newcommand*\listoflistings{\listof{codelisting}{List of Listings}}
\makeatother
\makeatletter
\makeatother
\makeatletter
\@ifpackageloaded{caption}{}{\usepackage{caption}}
\@ifpackageloaded{subcaption}{}{\usepackage{subcaption}}
\makeatother
\usepackage{bookmark}
\IfFileExists{xurl.sty}{\usepackage{xurl}}{} % add URL line breaks if available
\urlstyle{same}
\hypersetup{
  pdftitle={Energy, Projection, and Constraint in Three-Phase Linear Networks},
  pdfauthor={David et al.},
  colorlinks=true,
  linkcolor={blue},
  filecolor={Maroon},
  citecolor={Blue},
  urlcolor={Blue},
  pdfcreator={LaTeX via pandoc}}


\title{Energy, Projection, and Constraint in Three-Phase Linear
Networks}
\author{David et al.}
\date{2026-02-19}
\begin{document}
\maketitle

\renewcommand*\contentsname{Table of contents}
{
\hypersetup{linkcolor=}
\setcounter{tocdepth}{3}
\tableofcontents
}

\section{Abstract}\label{abstract}

We present a rigorously specified classical-field model that separates
three claims often conflated in foundational discussions of quantum
measurement. First, we prove that the Born-form dependence
\(\cos^2(\phi-\theta)\) arises exactly as a normalized energy projection
in a rotated mode basis of a linear two-dimensional internal space.
Second, we show in a transmission-line realization that a measurement
intervention is naturally represented as a boundary-condition change,
which instantaneously changes the admissible global modal decomposition
while leaving all physical propagation causal and finite-speed. Third,
we evaluate Clauser-Horne-Shimony-Holt (CHSH) correlations under
explicit detector models and verify numerically that local deterministic
dynamics with all trials counted satisfy \(S\le 2\), whereas
setting-dependent post-selection can produce apparent violations. The
contribution is therefore not a hidden-variable replacement for quantum
theory, but a transparent and auditable classical analog that sharply
distinguishes geometric projection, boundary-constraint update, and
genuinely nonclassical structure.

\section{1. Introduction}\label{introduction}

The main objective of this paper is to construct a single, internally
consistent framework in which:

\begin{enumerate}
\def\labelenumi{\arabic{enumi}.}
\tightlist
\item
  projection geometry and energy extraction are mathematically explicit,
\item
  measurement interventions are represented by physical boundary
  constraints,
\item
  causality is enforced by hyperbolic dynamics, and
\item
  Bell-inequality assumptions are tracked at the level of
  implementation.
\end{enumerate}

This yields a precise boundary between what can be captured by classical
linear-wave models and what cannot.

\subsection{1.1 Scope and claims}\label{scope-and-claims}

The model captures:

\begin{enumerate}
\def\labelenumi{\arabic{enumi}.}
\tightlist
\item
  exact Born-form projection weights in a rotated internal basis,
\item
  collapse-like update as a change in admissible mode space,
\item
  strict finite-speed signal and energy propagation.
\end{enumerate}

The model does not capture:

\begin{enumerate}
\def\labelenumi{\arabic{enumi}.}
\tightlist
\item
  tensor-product Hilbert-space nonlocality,
\item
  general noncommutative operator algebras,
\item
  loophole-free Bell violations under local realism.
\end{enumerate}

\section{2. Theoretical framework}\label{theoretical-framework}

\subsection{2.1 Internal mode space and energy inner
product}\label{internal-mode-space-and-energy-inner-product}

Let the internal two-mode field be

\[
\mathbf v(t)=\begin{pmatrix}v_\alpha(t)\\ v_\beta(t)\end{pmatrix},
\qquad
\mathbf v(t)=\Re\{\mathbf c e^{i\omega t}\},
\]

with \(\mathbf c\in\mathbb C^2\). Define the time-averaged quadratic
form

\[
\langle \mathbf v,\mathbf w\rangle_T=\frac1T\int_0^T \mathbf v(t)^\top\mathbf w(t)\,dt.
\]

For narrowband steady-state signals,

\[
\langle \mathbf v,\mathbf v\rangle_T\propto \|\mathbf c\|_2^2,
\]

so the internal space inherits a Euclidean inner-product structure.

\subsection{2.2 Rotated analyzer and projection
theorem}\label{rotated-analyzer-and-projection-theorem}

Let the analyzer angle be \(\theta\) with rotation

\[
R(\theta)=
\begin{pmatrix}
\cos\theta & -\sin\theta\\
\sin\theta & \cos\theta
\end{pmatrix},
\qquad
\mathbf v' = R(-\theta)\mathbf v.
\]

Assume a dissipative element extracts energy from the \(\alpha'\)
channel only. If the internal phasor direction is \(\phi\), then the
extracted energy fraction is

\[
W_\theta
= \frac{\int (v'_{\alpha})^2 dt}{\int \|\mathbf v\|_2^2 dt}
=\cos^2(\phi-\theta).
\]

This is algebraically identical to the Born-form squared projection
amplitude in a real two-dimensional state space.

\subsection{2.3 Three-phase embedding via Clarke
coordinates}\label{three-phase-embedding-via-clarke-coordinates}

For balanced three-phase voltages \((v_a,v_b,v_c)^\top\), define the
power-invariant Clarke map

\[
\begin{pmatrix}v_\alpha\\v_\beta\\v_0\end{pmatrix}
=
\underbrace{\begin{pmatrix}
\sqrt{\frac23} & -\frac1{\sqrt6} & -\frac1{\sqrt6}\\
0 & \frac1{\sqrt2} & -\frac1{\sqrt2}\\
\frac1{\sqrt3} & \frac1{\sqrt3} & \frac1{\sqrt3}
\end{pmatrix}}_{T_C}
\begin{pmatrix}v_a\\v_b\\v_c\end{pmatrix}.
\]

In balanced operation \(v_0\approx 0\), so the relevant internal space
is the \(\alpha\beta\) plane.

\section{3. Measurement as boundary constraint in a causal wave
system}\label{measurement-as-boundary-constraint-in-a-causal-wave-system}

\subsection{3.1 Transmission-line
dynamics}\label{transmission-line-dynamics}

Let \(\mathbf v(x,t),\mathbf i(x,t)\in\mathbb R^3\) satisfy

\[
\partial_x\mathbf v = -L'\partial_t\mathbf i,
\qquad
\partial_x\mathbf i = -C'\partial_t\mathbf v,
\]

with \(L',C'>0\) and propagation speed \(c=1/\sqrt{L'C'}\).

Define energy density and flux

\[
u=\frac12 C'\|\mathbf v\|_2^2+\frac12 L'\|\mathbf i\|_2^2,
\qquad
S=\mathbf v^\top\mathbf i,
\]

so \(\partial_t u+\partial_x S=0\).

\subsection{3.2 Detector boundary
condition}\label{detector-boundary-condition}

At a boundary port, the detector is a conductance matrix
\(G_{abc}(\theta)\succeq0\):

\[
\mathbf i(x_{\mathrm{port}},t)=G_{abc}(\theta)\,\mathbf v(x_{\mathrm{port}},t).
\]

Absorbed power is nonnegative:

\[
P_{\mathrm{abs}}(t)=\mathbf v^\top G_{abc}(\theta)\mathbf v\ge 0.
\]

A measurement action is modeled as switching \(G_{abc}\) at time
\(t_0\), i.e.~a change in boundary constraints.

\subsection{3.3 Domain update versus signal
propagation}\label{domain-update-versus-signal-propagation}

Let \(\mathcal D_{\mathrm{old}}\) and \(\mathcal D_{\mathrm{new}}\)
denote field domains satisfying old and new boundary conditions. At
\(t_0^+\), the admissible modal decomposition changes from
\(\mathcal D_{\mathrm{old}}\) to \(\mathcal D_{\mathrm{new}}\). This
update is instantaneous at the level of description (function space
constraints), but does not imply superluminal transfer of energy or
information. Physical changes at distance \(d\) require at least
\(d/c\).

\section{4. Reproducible computational
experiments}\label{reproducible-computational-experiments}

\subsection{4.1 Figure 1: projection geometry in internal mode
space}\label{figure-1-projection-geometry-in-internal-mode-space}

\pandocbounded{\includegraphics[keepaspectratio]{paper_files/figure-pdf/cell-3-output-1.png}}

The extracted fraction is
\(\|\mathrm{proj}_\theta(\mathbf u)\|_2^2=\cos^2(\phi-\theta)\).

\subsection{4.2 Figure 2: Born-form energy
law}\label{figure-2-born-form-energy-law}

\pandocbounded{\includegraphics[keepaspectratio]{paper_files/figure-pdf/cell-4-output-1.png}}

\subsection{4.3 Figure 3: boundary switch and finite-speed
causality}\label{figure-3-boundary-switch-and-finite-speed-causality}

\pandocbounded{\includegraphics[keepaspectratio]{paper_files/figure-pdf/cell-5-output-1.png}}

\begin{verbatim}
First detectable right-end deviation: t = 1.972
Causal lower bound (t0 + L/c):       t = 2.000
\end{verbatim}

The observed first deviation should not precede \(t_0+L/c\); in this run
it is delayed until approximately that bound (within discretization
tolerance).

\subsection{4.4 Figure 4: CHSH audit under explicit counting
rules}\label{figure-4-chsh-audit-under-explicit-counting-rules}

\pandocbounded{\includegraphics[keepaspectratio]{paper_files/figure-pdf/cell-6-output-1.png}}

\begin{verbatim}
S (all trials counted):        2.000
S (post-selected sample only): 4.000
Mean coincidence efficiency in post-selected run: 0.340
\end{verbatim}

\section{5. Results}\label{results}

\subsection{5.1 Born-form identity}\label{born-form-identity}

Figure 2 confirms the exact law \(W_\theta=\cos^2(\phi-\theta)\),
derived directly from Euclidean projection in the rotated internal
basis. No stochastic postulate is required for this identity.

\subsection{5.2 Boundary-switch
causality}\label{boundary-switch-causality}

Figure 3 demonstrates two separate facts:

\begin{enumerate}
\def\labelenumi{\arabic{enumi}.}
\tightlist
\item
  the boundary change is applied at \(t=t_0\);
\item
  a remote observable changes only after at least \(L/c\).
\end{enumerate}

Thus the model supports instantaneous update of admissible modal
constraints while preserving strict finite-speed physical propagation.

\subsection{5.3 CHSH audit}\label{chsh-audit}

In the all-counted local deterministic run, the computed CHSH value
saturates the classical limit (\(S\approx2\) within Monte Carlo error),
as required by Bell assumptions. In the post-selected run with
setting-dependent efficiency, \(S\) can exceed 2 substantially while
coincidence efficiency falls, identifying explicit sampling bias rather
than nonlocal dynamics.

\section{6. Discussion}\label{discussion}

\subsection{6.1 What is reconstructed}\label{what-is-reconstructed}

The framework reconstructs three structural elements often associated
with quantum measurement:

\begin{enumerate}
\def\labelenumi{\arabic{enumi}.}
\tightlist
\item
  squared-projection weighting,
\item
  basis dependence of outcomes,
\item
  global-state reinterpretation under measurement constraints.
\end{enumerate}

\subsection{6.2 What remains genuinely
nonclassical}\label{what-remains-genuinely-nonclassical}

The construction does not reproduce the full architecture of quantum
theory. In particular, it does not produce loophole-free Bell violations
under locality and deterministic hidden variables, and it does not
provide a replacement for noncommutative operator structure on
tensor-product Hilbert spaces.

\subsection{6.3 Methodological value}\label{methodological-value}

The practical value is diagnostic: by making propagation, detector
physics, and event-selection rules explicit, the model prevents
conceptual slippage between geometric projection effects,
boundary-constraint updates, and claims of nonlocal physical signaling.

\section{7. Limitations and future
work}\label{limitations-and-future-work}

\begin{enumerate}
\def\labelenumi{\arabic{enumi}.}
\tightlist
\item
  The current numerical experiment uses a scalar TL envelope for the
  causality figure; extending to fully coupled \(3\times3\)
  per-unit-length matrices \((\mathbf L,\mathbf C)\) is straightforward
  and should be done for higher-fidelity engineering correspondence.
\item
  Detector electronics are represented by linear shunt conductances and
  a stylized threshold coincidence rule. Future work should include
  transient front-end dynamics and noise models calibrated to hardware.
\item
  The CHSH loophole figure is intentionally didactic. A full
  detection-loophole phase map \(S(\eta,\Delta,T)\) with confidence
  intervals should be reported in a dedicated simulation paper.
\end{enumerate}

\section{8. Conclusion}\label{conclusion}

This paper establishes a rigorous classical-field bridge with three
precise outputs:

\begin{enumerate}
\def\labelenumi{\arabic{enumi}.}
\tightlist
\item
  Born-form dependence emerges exactly as normalized projection energy
  in a rotated internal mode basis.
\item
  Collapse-like update is formally represented as a boundary-domain
  change in admissible global modes.
\item
  Causality is preserved at the dynamical level, and Bell limits remain
  intact under local all-counted assumptions.
\end{enumerate}

The framework is therefore best interpreted as a structural analog and
clarification tool, not as a reduction of quantum mechanics to classical
circuitry.

\section{Appendix A: Minimal derivation of the projection
law}\label{appendix-a-minimal-derivation-of-the-projection-law}

Let a unit internal vector be \(\mathbf u=(\cos\phi,\sin\phi)^\top\) and
analyzer axis \(\hat e_\theta=(\cos\theta,\sin\theta)^\top\). Then

\[
\mathrm{proj}_{\theta}(\mathbf u)=(\mathbf u\cdot\hat e_\theta)\hat e_\theta,
\]

and

\[
\|\mathrm{proj}_{\theta}(\mathbf u)\|_2^2=(\mathbf u\cdot\hat e_\theta)^2
=(\cos\phi\cos\theta+\sin\phi\sin\theta)^2
=\cos^2(\phi-\theta).
\]

\section{Appendix B: Reproducibility
notes}\label{appendix-b-reproducibility-notes}

\begin{enumerate}
\def\labelenumi{\arabic{enumi}.}
\tightlist
\item
  All figures are generated in-document by Python code blocks.
\item
  Random sampling is seeded
  (\texttt{rng\ =\ np.random.default\_rng(20260219)}).
\item
  The CHSH example is intentionally split into all-counted and
  post-selected regimes to keep assumptions explicit.
\end{enumerate}




\end{document}
